\chapter{Relief: The Shaping of Ground}
\label{ch:relief}

\section{Relief as a Solve}

Relief is how flat ground becomes terrain. It is stated per group, and it is a
\emph{solve}: the studio takes a small number of statements about where the ground
must be, and computes a height for every cell in between.

The single most important property follows from that:

\begin{quote}
A relief replaces the top of every column it solves, and the heights it solves come
from its marks and pushes. A solved shape's own \id{base\_height} therefore decides
nothing about where the ground ends up.
\end{quote}

The technique cards prove this rather than assert it. The same relief is applied to
two plots whose shapes are drawn at different heights, 13 and 9 in one and both 9
in the other, and the built ground is identical: 2304 columns compared, 0 differ. An author who tunes a shape's drawn height to influence a solved surface is
tuning a number nothing reads.

\section{Constituents of a Relief}

\begin{lstlisting}[language=json]
"relief": { "stepped": {
  "base": 11, "reach": 0, "step": 1,
  "marks": [
    { "id": "upland-bench", "kind": "area", "h": 13, "bevel": 2,
      "ring": [[-28,-26],[-20,-30],[-8,-29],[-3,-22],
               [-4,-12],[-10,-7],[-22,-8],[-28,-14]] },
    { "id": "lowland-bench", "kind": "area", "h": 9, "bevel": 2,
      "ring": [[-46,8],[-38,5],[-20,6],[-6,5],
               [-2,10],[-8,15],[-26,14],[-44,15]] } ],
  "pushes": [
    { "id": "knoll", "amount": 1, "falloff": 3, "crown": 3.5,
      "roughness": 0, "seed": 1,
      "ring": [[-32,-32],[-21,-32],[-18,-25],[-22,-18],[-32,-19]] } ] } }
\end{lstlisting}

A \textbf{mark} pins ground. An \id{area} mark holds the middle of a region at a
stated height, and its \id{bevel} softens the edge. A \textbf{push} raises a region
by an amount over a falloff, with a crown and an optional roughness.

\Cref{tab:marks} is the full set. A mark is placed rather than sampled, for the
same reason a prop is: a decision about where the high ground is decides how the map
plays, so it belongs to the person making it rather than to a noise field.

\begin{table}[htbp]
\caption{The five marks. Each pins part of the footprint at a chosen height; the
solve fills between them.}
\label{tab:marks}
\small
\begin{tabularx}{\textwidth}{L{1.7cm} L{5.2cm} Y}
\toprule
\textbf{Mark} & \textbf{Pins} & \textbf{Says} \\
\midrule
\id{point} & a disc of a given radius & a summit, a hollow, a spot height \\
\id{line} & a band reaching \id{r} either side of a polyline, optionally with
  several heights along its run and a tread held flat &
  a ridge, a valley floor, a shoulder falling as it runs, a road switching back
  down a hill \\
\id{area} & every cell inside a ring, at one height or one per vertex, with an
  optional graded bevel inside its edge &
  a bench, a mesa top, a sunken floor, a shelf cut into a hillside \\
\id{rim} & the footprint's own outer rings & where the land meets the void \\
\id{scarp} & a band either side of a drawn line, at two heights, with the face
  between them left free &
  a break of slope, and the mark that decides where players can go \\
\bottomrule
\end{tabularx}
\end{table}

A line mark's \id{r} reaches either side of the centreline, so the band it writes is
\emph{twice} it: a stated 6 holds a twelve-block strip. The rim is optional, and it
is what keeps a hill inside its shape. Without one, marks alone decide the whole
surface, so a shape carrying a single high mark rises to that height everywhere and
runs off its own edge.

The most useful sentence about marks is about the space \emph{between} them:

\begin{quote}
Each \id{area} mark holds the middle of one meadow at its own height, and its
\id{bevel} softens the edge. The ground between the two benches is pinned by
nothing, so the solve grades it. \textbf{That free band is the connection.}
\end{quote}

Relief is therefore authored by pinning as little as possible. Pin the two things
that must be at particular heights and let the solve invent the slope between them;
pin everything and there is no slope, only a set of plateaux.

\begin{ruling}{Agreement of a Push's Gradients}
A push publishes two gradients: its skirt, which is \id{amount / falloff}, and its
crown. A knoll reading 0.33 up its skirt and 0.31 to its crown is smooth. Where the
two disagree the push steps at its own outline instead of rising into the ground,
and the studio raises \id{RL6}.
\end{ruling}

\section{Relief Scope: Three Treatments}

The question is what a relief does to a shape that already has a height. There are
three answers and the difference between two of them is the difference between a
join and a cliff.

\begin{description}[leftmargin=1.6em,style=nextline,itemsep=5pt]
  \item[No scope at all. The shape joins the solve.] Its drawn height is
    discarded and the relief decides where its ground goes. This is the default and
    it is right for terrain.
  \item[\id{hold}. The shape keeps its height and is still arrived at.] Its cells
    stay in the solve but are pinned at the height the shape states, so the ground
    around it is solved \emph{knowing where it has to arrive}.
  \item[\id{exclude}. The shape is a hole in the solve.] The relief bends round
    its edge exactly as it bends round the void, and nothing climbs to meet it. The
    shape keeps its faces.
\end{description}

\begin{figure}[htbp]
\centering
\includegraphics[width=\textwidth]{techniques/relief-on-shapes.png}

\medskip
{\small\begin{minipage}{0.95\textwidth}
The same three shapes, built three times under three treatments, sixty-four blocks
apart in one world. Each is a raised works yard, an upland meadow beside it and a
lowland below both. \textbf{Left}, the relief smooths the two meadows into one slope
and runs up against the yard without moving it. \textbf{Centre}, both meadows are
drawn at one height and the same relief builds it alike, block for block.
\textbf{Right}, the upland is held at the height it was drawn at and the relief is
spent only on the grade from the lowland up to it.
\end{minipage}}
\caption{The three treatments, side by side rather than in three folders. The
comparison is controlled: the shapes differ only in the heights they are drawn at
and in what each states about its scope.}
\label{fig:relief-scope}
\end{figure}

The distinction between the last two is measurable, and the technique card measures
it. Along one line, a shape drawn at 12 dropping to ground drawn at 8:

\begin{table}[htbp]
\caption{The same shape under \id{hold} and under \id{exclude}, read along one
line. \id{hold} makes the join; \id{exclude} leaves a face.}
\label{tab:hold-exclude}
\small
\begin{tabular}{l r r r r r r r r r}
\toprule
$z$ & $\cdots-1$ & 0 & 1 & 2 & 3 & 4 & 5 & 6 & 7 \\
\midrule
as drawn                        & 12 & 8 & 8 & 8 & 8 & 8 & 8 & 8 & 8 \\
\id{hold}, with the relief      & 12 & 12 & 11 & 11 & 10 & 10 & 9 & 9 & 8 \\
\id{exclude} instead            & 12 & 8 & 8 & 8 & 8 & 8 & 8 & 8 & 8 \\
\bottomrule
\end{tabular}
\end{table}

\begin{quote}
\id{exclude} on the upland would not make the join. An excluded shape is a
\textbf{hole in the solve}, the relief bends round its edge as it bends round the
void, and nothing climbs to meet it.
\end{quote}

The corresponding failure runs the other way. A built thing, such as a yard, a
plinth or a platform, wants \id{exclude}, because it should keep its faces. The same
layout built without it puts the yard's top at 11 instead of 17: the relief simply
replaced the height it was drawn at.

\section{Joins Between Grounds}

This is the craft rule the whole chapter exists to support, and it is stated in the
authoring discipline rather than in the API.

A board is two or three or four \emph{grounds}, each stated by the one instrument
that can state it, and every join between them is chosen rather than left over.
A beach and the dunes behind it are two reliefs on two layers, held at different
heights, because relief solves per layer and the fields are shifted into world $Y$
separately, so two reliefs can meet. One relief across both is one field, and
there is no join at all.

That is the difference between a board that reads as a place and a board that
reads as one field with a gradient in it.

\begin{table}[htbp]
\caption{A shore board, worked through: each ground stated by the instrument that
can state it, and each join chosen.}
\label{tab:shore}
\small
\begin{tabularx}{\textwidth}{L{3.1cm} L{4.1cm} Y}
\toprule
\textbf{The ground} & \textbf{Stated by} & \textbf{Why that instrument} \\
\midrule
the beach & a relief with flow: marks pinning only the waterline, a grain, a low
  step & it has to read as deposited, so nothing is pinned that does not have to be \\
the dunes behind it & a second relief, on its own layer, held a little higher &
  two reliefs can meet; one relief across both is one field and no join \\
the yard the houses stand on & a shape carrying \id{relief\_scope: "exclude"} &
  the footprint leaves the solve, so the two tiers meet at a \emph{face} \\
beach $\rightarrow$ dune & the two reliefs' own meeting, across the layer seam &
  the join has a shape because two fields made it \\
dune $\rightarrow$ yard & an authored flight: \id{height\_mode: "level"},
  anchors, \id{skirt: 0}, a material rather than a theme, run at least twice the
  rise & a relief graded across the seam deletes the boundary; a flight states it \\
\bottomrule
\end{tabularx}
\end{table}

\begin{ruling}{One Map, Several Grounds}
Plains everywhere is not simplicity. It is one instrument.
\end{ruling}

\section{Shapes Laid on Solved Ground}

Sand on the beach, gravel on the tide line: something that paints rather than
shapes. Such a shape carries a theme, a \id{base\_height} equal to the ground it
lies on, and, in the part that goes wrong, \textbf{no} \id{relief\_scope} at all.

Two reasons. A shape owns the paint only on a cell it forms the surface of, so one
stated lower runs under the ground, paints nothing, and says so on a 200. And a
\id{relief\_scope} is a statement about \emph{height}: \id{follow} seats the shape
on the solved field and then pins its whole ring there rigidly, which over a
terrain patch is a floor. The strand solves level to the block, and every mark
inside it is overwritten in silence.

\section{The Relief Read}

One call answers what the marks did to \emph{each other}, which no render can
show. It is
\route{POST /api/map/\{slug\}/sketch/relief/read}, given the stored layout.

It answers, per group: the cell count, the low and high, the relief, walk and
scramble tiers, the faces, the cliffs, and four fields that are the reason to ask
it at all.

\begin{description}[leftmargin=1.6em,style=nextline,itemsep=4pt]
  \item[\id{silentMarks}] every mark that landed nowhere. A mark that pinned nothing
    is invisible in every picture.
  \item[\id{seams}] the pairs that meet on a step. A seam naming a \emph{shape}
    rather than a mark is a \id{relief\_scope} pinning ground the marks were meant
    to shape.
  \item[\id{level} and \id{largestField}] the two numbers that say whether the
    ground has a shape at all. Over about 0.45 and 0.13 respectively, it is a table
    with edges.
  \item[\id{landform}] what the solve thinks it made: \emph{rolling},
    \emph{plain}, and so on.
\end{description}

\begin{ruling}{Relief and Symmetry Folding}
Any symmetry mode folds the relief across the centre, a group's \id{mirrors: false}
does not change that, and a layout stating no \id{mirror\_mode} takes
\id{rot\_180}. A board meant to solve unfolded has to say \id{mirror\_mode: "none"}.
\end{ruling}

\section{The Limits of Relief}

The repository's author states the danger directly: relief should be used, and
solid structural areas matter. A board whose every square is rolling has no ground
a fight can happen on, and an agent reaching for relief to make a board interesting
will reliably overdo it.

The counterpart discipline is a count. Run over a finished spec: how many reliefs,
how many marks and of what kinds, how many pushes, how many shapes at each height
mode, how much made ground, how many polylines, how many made layers, how many
copied trees. A zero in that list is not a fault. Four zeros is a board that used
one instrument and called it terrain.
